\scanpage{145}
% archive foliation: 144

caractérise \(\mathcal M\) à équivalence près). De plus, on
\unsure{aime|vise} donner des conditions moyennant
lesquelles le foncteur \(\underline h\) est \uline{pleinement}
\uline{fidèle}.

On va supposer que la structure \deleted{pré}additive
sur \(\mathcal V\) peut s'enrichir en une structure
\(\unsure{k|Q}\)-linéaire, où \(Q\) est un \uline{corps} commutatif,
[NB. si \(Q=\mathbb Q\), est évidemment déterminée
de façon unique]. On suppose que les
\deleted{\(\unsure{\operatorname{Hom}}\)}
\deleted{\(\operatorname{Fl}(X,Y)\) sont des vectoriels de dim
finie sur \(k\), et} que pour \(\forall X\), \(\operatorname{Fl}(X,X)\) est
une \uline{algèbre} \inserted{\uline{de dim finie}}
\uline{semi-simple} sur \(\deleted{k}\). Nous allons
alors \deleted{\unsure{en}} construire une catégorie \(\mathcal M^{\unsure{\sharp|*|\natural}}\) de la
façon suivante. Un objet de \(\mathcal M^{\unsure{\sharp|*|\natural}}\) est
un couple \((X,\pi)\), où \(X\in\operatorname{Ob}\mathcal V\), \(\pi\in\operatorname{Fl}(X,X)\),
\(\pi^2=\pi\) (\(\pi\) est un \uline{projecteur} de \(X\)). Un
homomorphisme \((X,\pi)\longrightarrow(X',\pi')\) est
un hom \(f:X\longrightarrow X'\) tel que
\(\pi'f=f\pi\), \uline{modulo} les \(f\) tels que
\(f\pi=\pi'f=0\). On compose de façon
évidente. D'où une catégorie, \inserted{\(\mathcal M^{\unsure{\sharp|*|\natural}}\)}, manifestement additive,
et un foncteur additif

\margin{i.e. \unsure{on a un hom} d'anneaux
\(Q\mathbin{-}\operatorname{End}\operatorname{id}_{\mathcal V}\)}

\diagram{
\begin{tikzcd}[column sep=large,row sep=large]
X \arrow[r,no head,"\pi"] \arrow[d,"f"'] &
X \arrow[d,no head,"f"]\\
X' \arrow[r,no head,"\pi'"] \arrow[d,no head,"g"'] &
X' \arrow[d,no head,"g"]\\
X'' \arrow[r,no head,"\pi''"] & X''
\end{tikzcd}
}
